\section{Introduction}

Suspended-payload UAV transport allows aerial systems to move objects that
cannot be rigidly attached to the vehicle. The suspended load, however, makes
the task sensitive to coupled vehicle-load dynamics
\cite{Sreenath2013DifferentiallyFlatHybrid,Li2023AutoTrans}. Under strong wind,
small tracking errors, payload swing, and transient speed changes can interact
with the cable-suspended payload and produce qualitatively different outcomes
across repeated runs \cite{Son2020ObstacleAvoidanceSuspendedLoad,GuerreroSanchez2017SwingAttenuation}.
Robust transport in this setting therefore requires not only feasible motion
generation and tracking, but also control over how aggressively the planned
motion is executed.

Payload-aware planning, MPC, and low-level control provide the core stack for
generating and tracking feasible trajectories
\cite{Li2023AutoTrans,Son2020ObstacleAvoidanceSuspendedLoad,UrbinaBrito2021PredictivePayloadTransport,Lee2010GeometricTrackingSE3}.
Improving these components alone does not fully address the runtime execution
problem. A fixed execution setting can work for one mission protocol or
wind/task phase, yet become too aggressive, too conservative, or poorly timed in
another. This creates a gap between nominal trajectory feasibility and balanced
protocol-level robustness under repeated strong-wind evaluation.

We study execution governance as a stack-compatible layer on top of an
AutoTrans-like suspended-payload UAV stack. The governor modifies command
execution through speed scale and acceleration scale, so it can be attached
externally without rewriting the planner, payload MPC, or SO(3) controller.
External reference or command modulation has a long control precedent
\cite{Garone2017ReferenceCommandGovernorsSurvey,Garone2016ExplicitReferenceGovernor,Li2021ActionGovernor},
but the governor studied here is empirical and is evaluated by repeated-run
behavior rather than by formal safety certification. In the learned variant,
risk-conditioned signals choose more conservative or less conservative
execution scales. This frames runtime aggressiveness as an interface-level
adaptation problem rather than a planner/controller redesign.

The evaluation must also separate protocol effects. A single-goal mission
protocol (goal repeat = 1) tests one commanded mission and is closest to
nominal single-goal execution. A goal-reissue stress protocol
(goal repeat = 10) repeatedly republishes the goal and stresses reference
updates, post-arrival behavior, and repeated command handling. Stress-testing
and benchmark-design studies show that protocol choice affects what a score
means \cite{Koren2018AdaptiveStressTesting,Nalic2020StressTestingScenarioBasedADS,Muller2019MeasureTargetConfusion}.
Because these protocols expose different failure modes and method rankings,
collapsing them into a mixed-protocol aggregate would obscure the relevant
robustness trade-offs.

The completed strong-wind evaluation shows that no current learned variant is
strongest in both protocols. Wind-Level 0.85 is the single-goal specialist,
achieving 26/30 strict-valid runs under goal repeat = 1, while Fixed Scale
0.80 is the goal-reissue stress specialist, achieving 24/30 strict-valid runs
under goal repeat = 10. The learned / risk-conditioned governor Risk Adapter
v1 is the most balanced current method, with the highest mean valid count
(24.0/30), highest worst-protocol valid count (23/30), and lowest total regret
(2). Risk Adapter v2.1 remains a strong nominal / single-goal variant, but it
is weaker under stress and is not framed as the final balanced method.

The paper contributes:

\begin{itemize}
  \item A stack-compatible learned / risk-conditioned execution governor for
  suspended-payload UAV transport that adapts speed scale and acceleration
  scale without replacing the planner, payload MPC, or SO(3) controller.
  \item We define a protocol-split strong-wind evaluation that separates the
  single-goal mission protocol (goal repeat = 1) from the goal-reissue
  stress protocol (goal repeat = 10), avoiding a mixed-protocol
  aggregate as the main result.
  \item Balanced robustness and protocol regret metrics showing that
  Risk Adapter v1 is the most balanced current learned / risk-conditioned
  method, while Wind-Level 0.85 and Fixed Scale 0.80 are strong protocol
  specialists.
  \item A failure-mode-aware diagnostic workflow based on
  invalid-only failure groups and representative traces, clarifying why
  aggregate strict-valid counts alone are insufficient for interpreting
  robustness.
\end{itemize}
